\let\negmedspace\undefined
\let\negthickspace\undefined
\documentclass[journal,12pt,onecolumn]{IEEEtran}
\usepackage{cite}
\usepackage{amsmath,amssymb,amsfonts,amsthm}
\usepackage{algorithmic}
\usepackage{graphicx}
\usepackage{textcomp}
\usepackage{xcolor}
\usepackage{txfonts}
\usepackage{listings}
\usepackage{enumitem}
\usepackage{mathtools}
\usepackage{gensymb}
\usepackage{comment}
\usepackage[breaklinks=true]{hyperref}
\usepackage{tkz-euclide} 
\usepackage{gvv}                                        
%\def\inputGnumericTable{}                                 
\usepackage[latin1]{inputenc}     
\usepackage{xparse}
\usepackage{color}                                            
\usepackage{array}                                            
\usepackage{longtable}                                       
\usepackage{calc}                                             
\usepackage{multirow}
\usepackage{multicol}
\usepackage{hhline}                                           
\usepackage{ifthen}                                           
\usepackage{lscape}
\usepackage{tabularx}
\usepackage{array}
\usepackage{float}
\newtheorem{theorem}{Theorem}[section]
\newtheorem{problem}{Problem}
\newtheorem{proposition}{Proposition}[section]
\newtheorem{lemma}{Lemma}[section]
\newtheorem{corollary}[theorem]{Corollary}
\newtheorem{example}{Example}[section]
\newtheorem{definition}[problem]{Definition}
\newcommand{\BEQA}{\begin{eqnarray}}
\newcommand{\EEQA}{\end{eqnarray}}
\newcommand{\define}{\stackrel{\triangle}{=}}
\theoremstyle{remark}
\newtheorem{rem}{Remark}
% Marks the beginning of the document
\begin{document}
\title{4.4.18}
\author{ai25btech11022 - Narshitha}
\maketitle
\renewcommand{\thefigure}{\theenumi}
\renewcommand{\thetable}{\theenumi}
\textbf{Question}:\\
If the point $\vec{A}\brak{x,y},\vec{B}\brak{3,6} $and $\vec{C}\brak{-3,4}$ are collinear,show that $x-3y+15=0$\\
\textbf{Solution}:\\
As $\vec{A},\vec{B} $and $\vec{C}$ are collinear\\
\begin{align}
    rank{\myvec{\vec{A-B}&& \vec{B-C}}}^T =1\\
    \myvec{x-3 && y-6 \\x+3 && y-4}
\end{align}
$R_2\longrightarrow R_2 -\frac{x+3}{x-3},R_1$
\begin{align}
    \myvec{x-3 && y-6 \\0 && y-4-\brak{y-6}\frac{x+3}{x-3}}\\
    \implies  y-4-\brak{y-6}\frac{x+3}{x-3} =0 \\
    \therefore x-3y+15=0
\end{align}
\begin{figure}
    \centering
    \includegraphics[width=0.9\linewidth]{figs/fig1.png}
    \caption{ }
    \label{fig:placeholder}
\end{figure}
\end{document}
